\chapter{Rebuilding and extending this manual}\label{chap:reproduction}
\section{A standalone local build}
The manual is intentionally isolated. Its source, generated reference data,
build helper, and final PDF are all under \api{manual/}. It adds no CMake
target and requires no change to the project's other documentation.

From the repository root, run:
\begin{lstlisting}[style=plain]
sh manual/build.sh
\end{lstlisting}
The script uses an already installed C++20 compiler and \api{pdflatex}.
It neither installs packages nor downloads resources. It compiles
\api{pictures/generate_reference.cpp} against \api{../include}, checks and
regenerates the geometry data, runs LaTeX three times to resolve the contents
and references, then writes \api{manual/nurbspath-manual.pdf}. Temporary
compiler and TeX products remain under \api{manual/build/}.

The local TeX installation must already provide TikZ, PGFPlots, Latin Modern,
AMS mathematics, geometry, listings, hyperref, booktabs/tabularx/longtable,
enumitem, needspace, fancyhdr, microtype, bookmark, and tcolorbox. This edition
was built with locally available TeX Live 2026. TikZ pictures are compiled
inside the document, without shell escape or externalisation.

\section{Source layout}
\begin{tabularx}{\linewidth}{@{}lY@{}}
\toprule
Location & Purpose\\
\midrule
\api{main.tex} & Master document and optional author-file hooks.\\
\api{preamble.tex} & Typography, colours, TikZ styles, code listings, and page layout.\\
\api{sections/00-title.tex} & Cover and snapshot identification.\\
\api{sections/01-foundations.tex} & Worlds, units, parameter conventions, and first program.\\
\api{sections/02-entities.tex} & Positions, vectors, rays, round surfaces, planes, embedding.\\
\api{sections/03-creators.tex} & Value construction and shared-ownership factories.\\
\api{sections/04-splines.tex} & Ordinary/array NURBS definitions and curve algorithms.\\
\api{sections/05-algorithms.tex} & Numerical utilities, intersections, and distances.\\
\api{sections/06-graphics-io.tex} & SVG rendering, persistence, and advisory output.\\
\api{sections/07-reference.tex} & API catalogue introduction.\\
\api{sections/reference_api.tex} & Generated signature/brief catalogue.\\
\api{pictures/} & C++ reference generator, sampled data, and display-value macros.\\
\api{tools/generate_api.py} & Optional local Doxygen-to-LaTeX catalogue refresh.\\
\bottomrule
\end{tabularx}

\section{Your introductory and closing material}
The master already contains conditional hooks for two files:
\begin{lstlisting}[style=plain]
\InputIfFileExists{intro.tex}{}{}
...
\InputIfFileExists{external.tex}{}{}
\end{lstlisting}
Place \api{intro.tex} in the manual directory to insert your introductory
words after the cover and before the contents. Place \api{external.tex} there
to insert your closing material after the appendices. Neither file is created
as a placeholder, so there is nothing to overwrite when you import your own
files. The document builds while they are absent.

These are included fragments, not separate documents: omit
\verb|\documentclass|, \verb|\begin{document}|, and
\verb|\end{document}|. You can use the same TikZ styles, \verb|\api|,
\api{manualnote}, and \api{lstlisting} formatting as the rest of the book.
Use \verb|\chapter*{Introduction}| or a similar unnumbered heading if wanted;
add it to the contents explicitly with \verb|\addcontentsline| if appropriate.

\section{Refreshing the API snapshot}
A normal build consumes the retained LaTeX catalogue. To refresh it from
the current public headers using an already installed Doxygen and Python:
\begin{lstlisting}[style=plain]
cd manual
python3 tools/generate_api.py
sh build.sh
\end{lstlisting}
The extractor writes a machine-readable coverage inventory to
\api{pictures/api_inventory.json}; it excludes private members and detail
helpers and counts public callable overloads. No network call or package
installation is part of this process. The installed Doxygen 1.18 comment
lexer can occasionally terminate by signal in this environment; the refresh
helper retries signal failures a bounded number of times and never suppresses
an ordinary documentation error.

\section{Checking numerical and visual changes}
The C++ reference generator checks interpolation at the supplied stations,
rational quarter-circle radii, crossing and tangent contact counts, closest
point consistency, array conversion, plane embedding, root discovery, and a
linear-system solution before writing coordinates. Display rounding occurs
after those checks. It is a reproducible worked-example validation, not a
replacement for the library's numbered test suite.

After editing the book, inspect its build log for LaTeX errors, unresolved
references, and overfull boxes. Render pages with a local PDF renderer to
check diagrams, tables, code, contents, headers, and page breaks. The supplied
PDF is the rendered result of the checked-in section files and reference data.
